\section{agent\_usage}

Endnutzer-Workflow: \textbf{Agenten} w\textbackslash\{\}"ahlen, \textbf{Agent Tasks} nutzen und \textbf{Delegation} \textbackslash\{\}"uber \texttt{/delegate}. Codebasis: \texttt{app/agents/}, Workspaces \texttt{operations\_agent\_tasks}, \texttt{cc\_agents}, Slash-Parser \texttt{app/core/commands/chat\_commands.py}.

\subsection{Ziel}

Sie nutzen ein \textbf{Agentenprofil} oder lassen eine \textbf{Delegations}-Bearbeitung mit mehreren Schritten \textbackslash\{\}"uber die Agenten-Infrastruktur laufen.

\subsection{Voraussetzungen}

\begin{enumerate}
\item Anwendung gestartet, \textbf{Ollama} erreichbar (wie beim Chat).
\item Agenten sind in Ihrer Umgebung \textbf{angelegt} (Standard: Seed-Profile beim ersten Start --- siehe \texttt{app/agents/seed\_agents.py}; bei leerem Dropdown Admin informieren).
\item Sie kennen den Unterschied: \textbf{Agent im Chat} (andere Systemrolle) vs. \textbf{`/delegate`} (Orchestrierungspfad).
\end{enumerate}

\subsection{Schritte A: Agent im Chat verwenden}

\begin{enumerate}
\item \textbf{Operations \$\textbackslash\{\}rightarrow\$ Chat} \textbackslash\{\}"offnen (siehe \textbf{chat\_usage}).
\item In der \textbf{Kopfzeile} des Chats (oder dem daf\textbackslash\{\}"ur vorgesehenen Steuerelement): \textbf{Agent} aus der Liste w\textbackslash\{\}"ahlen --- falls die Oberfl\textbackslash\{\}"ache das anbietet.
\item Nachricht normal eingeben und senden.
\item Das Modell arbeitet mit dem \textbf{Systemprompt} und der \textbf{Modellzuweisung} des gew\textbackslash\{\}"ahlten Profils (\texttt{AgentProfile}).
\end{enumerate}

\subsection{Schritte B: Agent Tasks \textbackslash\{\}"offnen}

\begin{enumerate}
\item Sidebar: \textbf{Operations}.
\item In der zweiten Leiste \textbf{Agent Tasks} w\textbackslash\{\}"ahlen (Workspace \textbf{`operations\_agent\_tasks`}).
\item Dort die projektbezogenen Aufgaben und Status einsehen und Aktionen ausf\textbackslash\{\}"uhren, die Ihre Version der Oberfl\textbackslash\{\}"ache anbietet (Listen, Buttons --- ohne hier Einzelfunktionen zu erfinden).
\end{enumerate}

\subsection{Schritte C: Delegation mit \texttt{/delegate}}

\begin{enumerate}
\item \textbf{Operations \$\textbackslash\{\}rightarrow\$ Chat} \textbackslash\{\}"offnen.
\item In das Eingabefeld schreiben --- \textbf{eine Zeile}, beginnend mit:
\end{enumerate}

\begin{itemize}
\item \textbf{`/delegate `} (Slash, Wort delegate, \textbf{Leerzeichen})
\item direkt danach \textbf{Ihre Aufgabe in eigenen Worten}, z. B.:
\end{itemize}

     \texttt{/delegate Fasse die letzten drei \textbackslash\{\}"Anderungen im Projekt in drei Bulletpoints zusammen}
\begin{enumerate}
\item Senden.
\item \textbf{Ohne Text nach `/delegate`} zeigt die App nur einen \textbf{Hinweis zur korrekten Verwendung} --- es wird keine Delegationsaufgabe mit Inhalt gestartet (so implementiert in \texttt{parse\_slash\_command}).
\end{enumerate}

\subsection{Varianten}

\begin{center}
\begin{tabular}{ll}
\hline
Situation & Vorgehen \\
\hline
Schnelle Frage ohne Spezialagent & Keinen Agent w\textbackslash\{\}"ahlen; Standard-Chat. \\
Code-fokussiert & Agent ?Code\texttt{} o. \textbackslash\{\}"A. w\textbackslash\{\}"ahlen, falls vorhanden, oder \texttt{/code} (siehe \textbf{chat\_usage}). \\
Gro\textbackslash\{\}ss\{\}e, mehrteilige Aufgabe & \texttt{/delegate} mit klarer Aufgabenbeschreibung in \textbf{einem} Senden. \\
Agenten verwalten & \textbf{Control Center \$\textbackslash\{\}rightarrow\$ Agents} (\texttt{cc\_agents}) --- Profile ansehen/bearbeiten, soweit die UI erlaubt. \\
\hline
\end{tabular}
\end{center}

\subsection{Fehlerf\textbackslash\{\}"alle}

\begin{center}
\begin{tabular}{ll}
\hline
Was Sie sehen & Was Sie tun \\
\hline
Agent-Dropdown leer & App neu starten; \textbf{Control Center \$\textbackslash\{\}rightarrow\$ Agents} pr\textbackslash\{\}"ufen; Admin wegen Datenbank/Seed benachrichtigen. \\
\texttt{/delegate} ohne Reaktion auf die Aufgabe & Pr\textbackslash\{\}"ufen, ob \textbf{Text nach} \texttt{/delegate} stand; nur \texttt{/delegate} allein reicht nicht. \\
Fehlermeldung ?Unbekannter Befehl\texttt{} & Tippfehler: exakt \textbf{`/delegate`} schreiben. \\
Delegation l\textbackslash\{\}"auft, Ergebnis unbefriedigend & Aufgabe \textbf{sch\textbackslash\{\}"arfer} formulieren; k\textbackslash\{\}"urzeren Umfang pro Aufruf; anderes Modell/Agent testen. \\
Agent Tasks zeigen nichts & \textbf{Projekt} aktiv? Agent Tasks sind projektbezogen, wenn ein Projekt ausgew\textbackslash\{\}"ahlt ist (siehe \texttt{docs/00\_map\_of\_the\_system.md}, global vs. projektbezogen). \\
\hline
\end{tabular}
\end{center}

\subsection{Tipps}

\begin{itemize}
\item \textbf{Eine klare Aufgabe pro `/delegate`}: Wer soll was liefern, in welchem Format?
\item Vor \textbf{Delegation} kurz pr\textbackslash\{\}"ufen, ob \textbf{Chat-Kontext} passt (Workflow \textbf{context\_control}) --- sonst fehlt dem Modell der gew\textbackslash\{\}"unschte Bezug.
\item F\textbackslash\{\}"ur wiederkehrende Agenten-Einstellungen \textbf{Control Center \$\textbackslash\{\}rightarrow\$ Agents} mit dem Team abstimmen, nicht jede Session neu erfinden.
\end{itemize}
